%% ma: L12-B03-0013
%% lop: 12
%% chuong: 1
%% bai: 3
%% dang: 12.3.3
%% loai: TN4
%% muc: TH
%% dap_an: "B"
%% nguon: "(NBV)-ĐỀ SỐ 15-MỖI NGÀY 1 ĐỀ THI 2025.pdf (D15, MỖI NGÀY 1 ĐỀ - Đề số 15), Phần I Câu 5"
%% kien_thuc: "Bài 3 (tiệm cận xiên)"
%% trang_thai: chinh-thuc
%% ghi_chu: "Trùng ý tưởng với L12-B03-0004 (cùng dạng tiệm cận xiên, hỏi ngược: chọn hàm nhận y = x + 1 làm tiệm cận xiên). Đáp án gốc B đúng."
\begin{ex}
Trong các hàm số sau, hàm số nào nhận đường thẳng $y=x+1$ làm tiệm cận xiên?
\choice{$y=\dfrac{x^2+x+2}{x+1}$}{\True $y=\dfrac{x^2+2x+3}{x+1}$}{$y=\dfrac{-x^2-x+2}{-x-1}$}{$y=\dfrac{-x^2+2}{x+1}$}
\loigiai{$\dfrac{x^2+2x+3}{x+1}=x+1+\dfrac2{x+1}$ nên tiệm cận xiên là $y=x+1$. Các hàm còn lại có tiệm cận xiên $y=x$, $y=x$ và $y=-x+1$.}
\end{ex}
